\section{Backup e Recovery}

\subsection{Estrategia de Backup}

\subsubsection{Tipos de Backup}

\paragraph{Backup Completo Diario}
\begin{lstlisting}[caption=Script de backup completo]
#!/bin/bash
# backup_completo.sh

DB_NAME="finops"
BACKUP_DIR="/var/backups/finops/completo"
DATE=$(date +%Y%m%d_%H%M%S)
LOGFILE="/var/log/finops_backup.log"

echo "$(date): Iniciando backup completo" >> $LOGFILE

# Backup com compressao
pg_dump -h localhost -U finops_backup -d $DB_NAME \
    --format=custom --compress=9 \
    --file="$BACKUP_DIR/finops_full_$DATE.backup"

if [ $? -eq 0 ]; then
    echo "$(date): Backup completo concluido com sucesso" >> $LOGFILE
    
    # Testar integridade do backup
    pg_restore --list "$BACKUP_DIR/finops_full_$DATE.backup" > /dev/null
    
    if [ $? -eq 0 ]; then
        echo "$(date): Backup validado com sucesso" >> $LOGFILE
    else
        echo "$(date): ERRO: Backup corrompido" >> $LOGFILE
        exit 1
    fi
else
    echo "$(date): ERRO: Falha no backup completo" >> $LOGFILE
    exit 1
fi

# Limpeza de backups antigos (manter 7 dias)
find $BACKUP_DIR -name "finops_full_*.backup" -mtime +7 -delete
\end{lstlisting}

\paragraph{Backup Incremental}
\begin{lstlisting}[caption=Configuracao de WAL archiving]
-- No postgresql.conf
-- wal_level = replica
-- archive_mode = on
-- archive_command = 'cp %p /var/backups/finops/wal/%f'
-- max_wal_senders = 3
-- checkpoint_timeout = 5min
-- max_wal_size = 1GB
-- min_wal_size = 80MB

-- Script para backup incremental
#!/bin/bash
# backup_incremental.sh

PGDATA="/var/lib/postgresql/data"
BACKUP_DIR="/var/backups/finops/incremental"
DATE=$(date +%Y%m%d_%H%M%S)

# Iniciar backup base
pg_basebackup -h localhost -U finops_backup \
    -D "$BACKUP_DIR/base_$DATE" \
    --format=tar --gzip --progress \
    --checkpoint=fast
\end{lstlisting}

\subsection{Procedures de Recovery}

\subsubsection{Recovery Completo}

\begin{lstlisting}[caption=Script de restore completo]
#!/bin/bash
# restore_completo.sh

if [ -z "$1" ]; then
    echo "Uso: $0 <arquivo_backup>"
    exit 1
fi

BACKUP_FILE="$1"
DB_NAME="finops"
TEMP_DB="finops_restore_$(date +%s)"

echo "Iniciando restore do backup: $BACKUP_FILE"

# Criar banco temporario
createdb -U postgres "$TEMP_DB"

# Restaurar dados
pg_restore -h localhost -U postgres -d "$TEMP_DB" \
    --verbose --clean --if-exists "$BACKUP_FILE"

if [ $? -eq 0 ]; then
    echo "Restore concluido no banco temporario: $TEMP_DB"
    echo "Para aplicar:"
    echo "1. Pare a aplicacao"
    echo "2. Renomeie o banco atual: ALTER DATABASE $DB_NAME RENAME TO ${DB_NAME}_old;"
    echo "3. Renomeie o banco restaurado: ALTER DATABASE $TEMP_DB RENAME TO $DB_NAME;"
    echo "4. Reinicie a aplicacao"
else
    echo "ERRO no restore"
    dropdb -U postgres "$TEMP_DB"
    exit 1
fi
\end{lstlisting}

\subsubsection{Point-in-Time Recovery}

\begin{lstlisting}[caption=Recovery para um momento especifico]
#!/bin/bash
# pitr_recovery.sh

if [ -z "$1" ] || [ -z "$2" ]; then
    echo "Uso: $0 <backup_base> <timestamp>"
    echo "Exemplo: $0 /path/backup.tar '2024-01-15 14:30:00'"
    exit 1
fi

BACKUP_BASE="$1"
TARGET_TIME="$2"
RECOVERY_DIR="/var/lib/postgresql/recovery"
WAL_DIR="/var/backups/finops/wal"

# Parar PostgreSQL
systemctl stop postgresql

# Limpar diretorio de dados
rm -rf /var/lib/postgresql/data/*

# Extrair backup base
tar -xzf "$BACKUP_BASE" -C /var/lib/postgresql/data/

# Criar arquivo de recovery
cat > /var/lib/postgresql/data/recovery.conf << EOF
restore_command = 'cp $WAL_DIR/%f %p'
recovery_target_time = '$TARGET_TIME'
recovery_target_timeline = 'latest'
EOF

# Iniciar PostgreSQL em modo recovery
systemctl start postgresql

echo "Recovery iniciado para o momento: $TARGET_TIME"
echo "Monitore os logs: tail -f /var/log/postgresql/postgresql.log"
\end{lstlisting}

\subsection{Monitoramento de Backup}

\subsubsection{Verificacao Automatica}

\begin{lstlisting}[caption=Script de verificacao de backups]
#!/bin/bash
# verificar_backups.sh

BACKUP_DIR="/var/backups/finops"
LOGFILE="/var/log/finops_backup_check.log"
EMAIL_ADMIN="admin@finops.com"

echo "$(date): Iniciando verificacao de backups" >> $LOGFILE

# Verificar backup mais recente
ULTIMO_BACKUP=$(find $BACKUP_DIR/completo -name "finops_full_*.backup" -mtime -1 | head -1)

if [ -z "$ULTIMO_BACKUP" ]; then
    echo "$(date): ERRO: Nenhum backup encontrado nas ultimas 24h" >> $LOGFILE
    echo "ALERTA: Backup nao encontrado" | mail -s "FinOps Backup Alert" $EMAIL_ADMIN
    exit 1
fi

# Testar integridade
pg_restore --list "$ULTIMO_BACKUP" > /dev/null 2>&1

if [ $? -eq 0 ]; then
    echo "$(date): Backup integro: $ULTIMO_BACKUP" >> $LOGFILE
else
    echo "$(date): ERRO: Backup corrompido: $ULTIMO_BACKUP" >> $LOGFILE
    echo "ALERTA: Backup corrompido" | mail -s "FinOps Backup Alert" $EMAIL_ADMIN
    exit 1
fi

# Verificar espaco em disco
ESPACO_LIVRE=$(df $BACKUP_DIR | awk 'NR==2 {print $4}')
LIMITE_MINIMO=1048576  # 1GB em KB

if [ $ESPACO_LIVRE -lt $LIMITE_MINIMO ]; then
    echo "$(date): AVISO: Pouco espaco em disco: ${ESPACO_LIVRE}KB" >> $LOGFILE
    echo "AVISO: Pouco espaco para backups" | mail -s "FinOps Storage Alert" $EMAIL_ADMIN
fi

echo "$(date): Verificacao concluida com sucesso" >> $LOGFILE
\end{lstlisting}

\subsection{Disaster Recovery}

\subsubsection{Plano de Contingencia}

\begin{lstlisting}[caption=Procedimento de disaster recovery]
-- PLANO DE DISASTER RECOVERY FINOPS
-- ====================================

-- FASE 1: AVALIACAO DO DANO
-- 1. Verificar se servidor PostgreSQL responde
-- 2. Testar conectividade de rede
-- 3. Verificar integridade dos discos
-- 4. Avaliar logs de erro

-- FASE 2: RECOVERY DE EMERGENCIA
-- 1. Provisionar novo servidor se necessario
-- 2. Instalar PostgreSQL na mesma versao
-- 3. Configurar parametros basicos
-- 4. Restaurar backup mais recente

-- Script de recovery de emergencia
#!/bin/bash
# emergency_recovery.sh

echo "=== FinOps Emergency Recovery ==="
echo "Este script ira restaurar o banco de dados"
echo "Confirme antes de prosseguir!"
read -p "Continuar? (y/N): " confirm

if [ "$confirm" != "y" ]; then
    echo "Recovery cancelado"
    exit 0
fi

# Encontrar backup mais recente
ULTIMO_BACKUP=$(find /var/backups/finops/completo -name "finops_full_*.backup" | sort | tail -1)

if [ -z "$ULTIMO_BACKUP" ]; then
    echo "ERRO: Nenhum backup encontrado"
    exit 1
fi

echo "Usando backup: $ULTIMO_BACKUP"

# Recriar banco
dropdb --if-exists finops
createdb finops

# Restaurar dados
pg_restore -d finops "$ULTIMO_BACKUP"

echo "Recovery concluido"
echo "Verificar logs e testar aplicacao antes de colocar em producao"
\end{lstlisting}

\subsubsection{Replicacao para Alta Disponibilidade}

\begin{lstlisting}[caption=Configuracao de replica]
-- No servidor master (postgresql.conf)
-- wal_level = replica
-- max_wal_senders = 3
-- wal_keep_segments = 64
-- synchronous_standby_names = 'standby1'

-- No pg_hba.conf do master
-- host replication finops_repl standby_ip/32 md5

-- Script para configurar standby
#!/bin/bash
# setup_standby.sh

MASTER_HOST="master_ip"
STANDBY_DATA="/var/lib/postgresql/standby"

# Parar PostgreSQL no standby
systemctl stop postgresql

# Backup inicial do master
pg_basebackup -h $MASTER_HOST -U finops_repl \
    -D $STANDBY_DATA --write-recovery-conf \
    --checkpoint=fast

# Configurar recovery.conf
cat >> $STANDBY_DATA/recovery.conf << EOF
standby_mode = 'on'
primary_conninfo = 'host=$MASTER_HOST port=5432 user=finops_repl'
trigger_file = '/var/lib/postgresql/failover'
EOF

# Iniciar standby
systemctl start postgresql

echo "Standby configurado. Para failover, crie o arquivo:"
echo "touch /var/lib/postgresql/failover"
\end{lstlisting}
